\chapter{Optimal Transport の基礎理論}
\label{ch:seminar-01}

%% ============================================================
%%  0. 最適割当問題（§2.2 前半）
%% ============================================================
\section{最適割当問題}
\label{sec:sem-assignment}

$n$ 人の作業者と $n$ 個の仕事があり，
作業者 $i$ を仕事 $j$ に配属するコストが $C_{i,j}$ で与えられているとする．
作業者と仕事を1対1に対応させるとき，
コストを最小にする割当を求めたい．

\begin{definition}{最適割当問題}{sem-assignment}
  コスト行列 $\mathbf{C} \in \R^{n \times n}$ が与えられたとき，
  \textbf{最適割当問題}（optimal assignment problem）とは，
  \[
    \min_{\sigma \in \Perm(n)} \frac{1}{n} \sum_{i=1}^{n} C_{i, \sigma(i)}
  \]
  を解く置換 $\sigma \in \Perm(n)$ を求める問題である．
\end{definition}

\begin{example}{最適割当}{sem-assignment-cost-example}
  3人の作業者を3つの仕事場に配属する状況を考える．
  各作業者の自宅と仕事場の位置が異なるため，
  通勤コスト $C_{i,j}$（作業者 $i$ が仕事場 $j$ に通う負担）は
  組合せによって異なる：
  \[
    \mathbf{C} =
    \begin{pmatrix}
      5 & 2 & 8 \\
      7 & 4 & 1 \\
      3 & 9 & 6
    \end{pmatrix}.
  \]
  たとえば $C_{1,2} = 2$ は
  「作業者1が仕事場2に通うコスト」を意味する．

  $3! = 6$ 通りの割当すべてを図示する．
  各図で青丸（{\color{blue!60}$\bullet$}）が作業者，
  赤四角（{\color{red!60}$\blacksquare$}）が仕事場を表し，
  矢印上の数字は個別コスト $C_{i,\sigma(i)}$ である．

  \begin{center}
  \begin{tikzpicture}[
    wk/.style={circle, draw=blue!60, fill=blue!10,
               minimum size=10pt, inner sep=1pt, font=\tiny},
    tk/.style={rectangle, rounded corners=1.5pt,
               draw=red!60, fill=red!10,
               minimum size=10pt, inner sep=1pt, font=\tiny},
    arr/.style={-{Stealth[length=3pt]}, semithick, black!50},
    optarr/.style={-{Stealth[length=3pt]}, semithick, green!50!black},
  ]
    % ---- helper macro: draw 6 nodes at spread positions ----
    \newcommand{\plotnodes}[1]{
      \node[wk] (w1#1) at (0.5, 3)   {1};
      \node[wk] (w2#1) at (3.5, 0.3) {2};
      \node[wk] (w3#1) at (0, 0.6)   {3};
      \node[tk] (t1#1) at (1.8, 0.6) {1};
      \node[tk] (t2#1) at (1.8, 3)   {2};
      \node[tk] (t3#1) at (2.8, 0)   {3};
    }
    %% === Row 1 ===
    % (1,2,3): 1→1, 2→2, 3→3, cost=5
    \begin{scope}[xshift=0cm]
      \plotnodes{A}
      \draw[arr] (w1A)--(t1A); \draw[arr] (w2A)--(t2A); \draw[arr] (w3A)--(t3A);
      \node[font=\scriptsize] at (1.8, -0.7) {(a)\; $\sigma=(1,2,3)$};
      \node[font=\scriptsize] at (1.8, -1.2) {コスト $= 5$};
    \end{scope}
    % (1,3,2): 1→1, 2→3, 3→2, cost=5
    \begin{scope}[xshift=4.5cm]
      \plotnodes{B}
      \draw[arr] (w1B)--(t1B); \draw[arr] (w2B)--(t3B); \draw[arr] (w3B)--(t2B);
      \node[font=\scriptsize] at (1.8, -0.7) {(b)\; $\sigma=(1,3,2)$};
      \node[font=\scriptsize] at (1.8, -1.2) {コスト $= 5$};
    \end{scope}
    % (2,1,3): 1→2, 2→1, 3→3, cost=5
    \begin{scope}[xshift=9cm]
      \plotnodes{C}
      \draw[arr] (w1C)--(t2C); \draw[arr] (w2C)--(t1C); \draw[arr] (w3C)--(t3C);
      \node[font=\scriptsize] at (1.8, -0.7) {(c)\; $\sigma=(2,1,3)$};
      \node[font=\scriptsize] at (1.8, -1.2) {コスト $= 5$};
    \end{scope}
    %% === Row 2 ===
    % (2,3,1): 1→2, 2→3, 3→1, cost=2 *** optimal ***
    \begin{scope}[xshift=0cm, yshift=-5cm]
      \draw[green!50!black, thick, rounded corners=4pt]
        (-0.4, -1.5) rectangle (4.0, 3.5);
      \plotnodes{D}
      \draw[optarr] (w1D)--(t2D); \draw[optarr] (w2D)--(t3D); \draw[optarr] (w3D)--(t1D);
      \node[font=\scriptsize, green!40!black] at (1.8, -0.7) {(d)\; $\sigma=(2,3,1)$};
      \node[font=\scriptsize\bfseries, green!40!black] at (1.8, -1.2) {コスト $= 2$（最適）};
    \end{scope}
    % (3,1,2): 1→3, 2→1, 3→2, cost=8
    \begin{scope}[xshift=4.5cm, yshift=-5cm]
      \plotnodes{E}
      \draw[arr] (w1E)--(t3E); \draw[arr] (w2E)--(t1E); \draw[arr] (w3E)--(t2E);
      \node[font=\scriptsize] at (1.8, -0.7) {(e)\; $\sigma=(3,1,2)$};
      \node[font=\scriptsize] at (1.8, -1.2) {コスト $= 8$};
    \end{scope}
    % (3,2,1): 1→3, 2→2, 3→1, cost=5
    \begin{scope}[xshift=9cm, yshift=-5cm]
      \plotnodes{F}
      \draw[arr] (w1F)--(t3F); \draw[arr] (w2F)--(t2F); \draw[arr] (w3F)--(t1F);
      \node[font=\scriptsize] at (1.8, -0.7) {(f)\; $\sigma=(3,2,1)$};
      \node[font=\scriptsize] at (1.8, -1.2) {コスト $= 5$};
    \end{scope}
  \end{tikzpicture}
  \end{center}

  各割当のコストをまとめると，
  \begin{center}
  \begin{tabular}{ccc}
     & $\sigma$ & コスト $\frac{1}{n}\sum_i C_{i,\sigma(i)}$ \\\hline
    (a) & $(1,2,3)$ & $\frac{1}{3}(5 + 4 + 6) = 5$ \\
    (b) & $(1,3,2)$ & $\frac{1}{3}(5 + 1 + 9) = 5$ \\
    (c) & $(2,1,3)$ & $\frac{1}{3}(2 + 7 + 6) = 5$ \\
    (d) & $(2,3,1)$ & $\frac{1}{3}(2 + 1 + 3) = \mathbf{2}$ \\
    (e) & $(3,1,2)$ & $\frac{1}{3}(8 + 7 + 9) = 8$ \\
    (f) & $(3,2,1)$ & $\frac{1}{3}(8 + 4 + 3) = 5$ \\
  \end{tabular}
  \end{center}
  よって $\sigma = (2, 3, 1)$ が最小コスト $2$ を達成し，最適である．
\end{example}

この例では $3! = 6$ 通りの全列挙で最適解が求まったが，
一般には $|\Perm(n)| = n!$ は急激に増大し，
$n = 70$ でも $n! > 10^{100}$ となるため全列挙は現実的でない．
ハンガリアン法やオークションアルゴリズムなど，
$O(n^3)$ で解ける効率的な手法が知られている．

\begin{remark}{最適割当問題の解の一意性}{sem-uniqueness}
  最適割当問題の最適解は\textbf{一意とは限らない}．

  \textbf{例}：$n = 2$ とし，
  作業者 $x_1, x_2$ と仕事場 $y_1, y_2$ が
  辺の長さ $1$ の正方形の頂点に配置されているとする．
  $x_1, x_2$ を対角頂点に，$y_1, y_2$ を残りの対角頂点に置くと，
  すべての作業者--仕事場間の距離が等しくなる（$= 1$）．

  \begin{center}
  \begin{tikzpicture}[scale=1.5]
    % Square outline (light)
    \draw[gray!30] (0,0) -- (1,0) -- (1,1) -- (0,1) -- cycle;
    % Workers (diagonal)
    \fill[blue] (0, 0) circle (3pt) node[below left] {\small $x_1$};
    \fill[blue] (1, 1) circle (3pt) node[above right] {\small $x_2$};
    % Tasks (other diagonal)
    \fill[red]  (1, 0) circle (3pt) node[below right] {\small $y_1$};
    \fill[red]  (0, 1) circle (3pt) node[above left] {\small $y_2$};
    % Assignment 1: σ = (1,2): x1→y1, x2→y2 — solid
    \draw[->, thick] (0.1, 0.02) -- (0.88, 0.02);
    \draw[->, thick] (0.9, 0.98) -- (0.12, 0.98);
    % Assignment 2: σ = (2,1): x1→y2, x2→y1 — dashed
    \draw[->, thick, dashed, gray] (0.02, 0.12) -- (0.02, 0.88);
    \draw[->, thick, dashed, gray] (0.98, 0.88) -- (0.98, 0.12);
    % Labels
    \node[below] at (0.5, -0.4) {
      \small $\sigma\!=\!(1,2)$: コスト $= 1$，\quad
      \textcolor{gray}{$\sigma\!=\!(2,1)$: コスト $= 1$}};
    % Side length
    \node[right, font=\scriptsize] at (1.15, 0.5) {辺の長さ $1$};
  \end{tikzpicture}
  \end{center}

  このとき $\sigma = (1,2)$ と $\sigma = (2,1)$ は
  ともにコスト $1$ を与え，いずれも最適である．
  すべての距離が等しいため，どの割当を選んでもコストは変わらない．
\end{remark}

%% ============================================================
%%  1. Monge 問題（§2.2 後半）
%% ============================================================
\section{Monge 問題}
\label{sec:sem-monge}

最適割当問題では，同数の点を均等な重みで1対1に対応させていた．
これを，個数や重みが異なる場合や，1対多の写像にも一般化したい．

「ある場所にある砂山を，最小の労力で別の場所に所定の形状で積み上げるには，
各砂粒をどこに運べばよいか？」---
1781年に Monge が提唱したこの問題は，
これらを許すより一般的な枠組みを与える．

\begin{definition}{Monge 問題（離散版）}{sem-monge-discrete}
  離散測度
  $\alpha = \sum_{i=1}^n a_i \delta_{x_i}$，
  $\beta = \sum_{j=1}^m b_j \delta_{y_j}$
  とコスト関数 $c \colon \X \times \Y \to \R_+$ が与えられたとする．
  \textbf{離散 Monge 問題}とは，
  \[
    \min_{T}
    \left\{
    \sum_{i=1}^{n} c(x_i, T(x_i))
    \;\middle|\;
    T\pushforward \alpha = \beta
    \right\}
  \]
  を求める問題である．
\end{definition}

特に $n = m$ かつ $a_i = b_j = 1/n$ のとき，
$T\pushforward \alpha = \beta$ は $T$ が全単射であることを要求するから，
$C_{i,j} \defeq c(x_i, y_j)$ とおけば
Monge 問題は最適割当問題に帰着する．

離散版を一般の確率測度に拡張すると，次の連続版が得られる．

\begin{definition}{Monge 問題（連続版）}{sem-monge}
  確率測度 $\alpha \in \Mm_+^1(\X)$，$\beta \in \Mm_+^1(\Y)$ と
  コスト関数 $c \colon \X \times \Y \to \R_+$ に対して，
  \textbf{Monge 問題}とは，
  \[
    \inf_{T \colon \X \to \Y}
    \left\{
    \int_{\X} c(x, T(x)) \, \d\alpha(x)
    \;\middle|\;
    T\pushforward \alpha = \beta
    \right\}
  \]
  を求める問題である．
\end{definition}

%% --- Monge 問題のイメージ図 ---
\begin{figure}[htbp]
\centering
\begin{tikzpicture}[scale=1.0]
  % --- X（土台）と α ---
  \begin{scope}
    \fill[gray!15] (-0.3, -0.4) rectangle (3.3, 0);
    \draw[thick] (-0.3, 0) -- (3.3, 0);
    \node[below] at (1.5, -0.4) {$\X$};
    % α
    \fill[blue!15] plot[smooth, tension=0.7]
      coordinates {(0,0) (0.4,0.5) (0.8,1.2) (1.2,1.8) (1.6,1.5)
                   (2.0,1.0) (2.4,0.5) (2.8,0.2) (3.0,0)}
      -- cycle;
    \draw[thick, blue] plot[smooth, tension=0.7]
      coordinates {(0,0) (0.4,0.5) (0.8,1.2) (1.2,1.8) (1.6,1.5)
                   (2.0,1.0) (2.4,0.5) (2.8,0.2) (3.0,0)};
    \node[blue] at (-0.3, 1.4) {$\alpha$};
    % T^{-1}(A) のハイライト
    \fill[green!35, opacity=0.7]
      (1.8, 0) -- plot[smooth, tension=0.7]
      coordinates {(1.8, 1.15) (2.0, 1.0) (2.2, 0.7) (2.4, 0.5)}
      -- (2.4, 0) -- cycle;
    \draw[thick, green!50!black, dashed] (1.8, 0) -- (1.8, 1.15);
    \draw[thick, green!50!black, dashed] (2.4, 0) -- (2.4, 0.5);
    \node[green!40!black, font=\small] at (2.1, 1.7) {$T^{-1}(A)$};
    \draw[->, green!40!black] (2.1, 1.5) -- (2.1, 1.15);
  \end{scope}

  % --- 写像 T（1本の矢印）---
  \draw[->, very thick, gray!80] (3.5, 1.0) -- (6.3, 1.0)
    node[midway, above, font=\small] {$T$};
  \node[gray, font=\scriptsize, text width=3cm, align=center] at (4.9, 0.5)
    {各点 $x$ の全質量を\\$T(x)$ に送る};

  % --- Y（土台）と β ---
  \begin{scope}[xshift=6.5cm]
    \fill[gray!15] (-0.3, -0.4) rectangle (3.3, 0);
    \draw[thick] (-0.3, 0) -- (3.3, 0);
    \node[below] at (1.5, -0.4) {$\Y$};
    % β
    \fill[red!15] plot[smooth, tension=0.7]
      coordinates {(0,0) (0.3,0.2) (0.6,0.5) (1.0,1.2) (1.4,1.7)
                   (1.8,1.3) (2.2,0.6) (2.6,0.2) (3.0,0)}
      -- cycle;
    \draw[thick, red] plot[smooth, tension=0.7]
      coordinates {(0,0) (0.3,0.2) (0.6,0.5) (1.0,1.2) (1.4,1.7)
                   (1.8,1.3) (2.2,0.6) (2.6,0.2) (3.0,0)};
    \node[red] at (3.3, 1.2) {$\beta$};
    % A のハイライト
    \fill[orange!35, opacity=0.7]
      (0.8, 0) -- plot[smooth, tension=0.7]
      coordinates {(0.8, 0.9) (1.0, 1.2) (1.2, 1.55) (1.4, 1.7) (1.5, 1.6)}
      -- (1.5, 0) -- cycle;
    \draw[thick, orange!70!black, dashed] (0.8, 0) -- (0.8, 0.9);
    \draw[thick, orange!70!black, dashed] (1.5, 0) -- (1.5, 1.6);
    \node[orange!60!black, font=\small] at (1.15, 2.2) {$A$};
    \draw[->, orange!60!black] (1.15, 2.0) -- (1.15, 1.55);
  \end{scope}

  % --- 押し出し条件 ---
  \draw[<-, thick, purple!70, dashed] (2.1, -0.8) to[bend left=15]
    node[below, font=\small] {逆像 $T^{-1}$}
    (7.65, -0.8);
  \node[font=\small] at (4.8, -1.8)
    {$T\pushforward \alpha = \beta
    \;\Longleftrightarrow\;
    \alpha\bigl(T^{-1}(A)\bigr) = \beta(A)
    \quad (\forall\, A)$};
\end{tikzpicture}
\caption{Monge 問題の概念図．写像 $T$ が $\X$ 上の測度 $\alpha$ を $\Y$ 上の測度 $\beta$ に押し出す．}
\label{fig:sem-monge}
\end{figure}

\subsection{Monge 問題の本質的困難}

Monge の定式化には2つの深刻な問題がある：
\begin{enumerate}
  \item \textbf{実行可能集合の非凸性}：凸最適化の手法が使えない．
  \item \textbf{解の非存在}：実行可能な写像 $T$ 自体が存在しない場合がある．
\end{enumerate}
順に見ていく．1つ目は凸性に関わるため，
以下では $\Y$ をベクトル空間（例：$\R^d$）に限定する．

\begin{definition}{凸結合・凸集合}{sem-convex}
  ベクトル空間 $V$ の元 $v_1, v_2 \in V$ と $t \in [0,1]$ に対して，
  $t v_1 + (1-t) v_2$ を $v_1$ と $v_2$ の\textbf{凸結合}という．
  部分集合 $S \subseteq V$ が\textbf{凸}であるとは，
  任意の $v_1, v_2 \in S$ の凸結合が再び $S$ に属すること，すなわち
  \[
    v_1, v_2 \in S,\; t \in [0,1]
    \;\Longrightarrow\;
    t v_1 + (1-t) v_2 \in S
  \]
  が成り立つことをいう．
\end{definition}

\begin{remark}{写像空間への適用}{sem-map-convex}
  以下で実行可能集合 $\mathcal{T}(\alpha,\beta)$ の凸性を議論するために，
  それが属する空間のベクトル空間構造を確認する．
  $\Y$ がベクトル空間のとき，$\X$ 上の写像全体 $\Y^{\X}$ は
  点ごとの演算によりベクトル空間をなす．
  写像 $T_1, T_2 \colon \X \to \Y$ の凸結合はこの構造のもとで
  \[
    \bigl(t T_1 + (1-t)T_2\bigr)(x) = t\,T_1(x) + (1-t)\,T_2(x)
    \qquad (\forall\, x \in \X)
  \]
  と計算される．したがって $\Y^{\X}$ の部分集合の凸性を問うことができる．
\end{remark}

\subsubsection{困難1：実行可能集合の非凸性}

Monge 問題の実行可能集合
\[
  \mathcal{T}(\alpha, \beta)
  \defeq \left\{ T \colon \X \to \Y \;\middle|\;
  T \text{ は可測},\; T\pushforward \alpha = \beta \right\}
\]
は $\Y^{\X}$ の部分集合であるが，一般に\textbf{凸でない}．
その理由は，押し出し $T\pushforward \alpha$ が写像 $T$ に関して
\textbf{非線形}であることにある．
すなわち，$T_1, T_2 \in \mathcal{T}(\alpha,\beta)$ のとき
${T_i}\pushforward \alpha = \beta$ であるから
$t \cdot {T_1}\pushforward \alpha + (1-t) \cdot {T_2}\pushforward \alpha
= t\beta + (1-t)\beta = \beta$
が成り立つが，凸結合の押し出しはこれと一致するとは限らない：
\[
  \bigl(t T_1 + (1-t)T_2\bigr)\pushforward \alpha
  \neq t \cdot {T_1}\pushforward \alpha + (1-t) \cdot {T_2}\pushforward \alpha
  \quad (\text{一般には}).
\]
直観的には，押し出しは各点の\textbf{行き先の分布}を見る操作であり，
行き先を点ごとに補間すると分布自体が変わってしまうためである．

\begin{remark}{測地線凸性}{sem-geodesic-convexity}
  ここでは $\Y$ がベクトル空間の場合に限り凸性を議論した．
  $\Y$ が一般の距離空間の場合，凸結合の代わりに2点間の測地線を用いた
  \textbf{測地線凸性}（geodesic convexity）で同様の議論が展開できる．
\end{remark}

\subsubsection{困難2：解の非存在}

写像による輸送では，質量の\textbf{分割}が許されない．
すなわち，$x$ の質量は $T(x)$ という1点にすべて送らなければならない．
このため，輸送元の測度が輸送先の測度より「粗い」場合に問題が生じる．

\begin{example}{Monge 写像が存在しない場合}{sem-monge-noexist}
  $\alpha = \delta_0$（原点に全質量を集中），
  $\beta = \frac{1}{2}\delta_{-1} + \frac{1}{2}\delta_1$
  （$-1$ と $1$ に半分ずつ配置）を考える．

  $T\pushforward \delta_0 = \delta_{T(0)}$ は1点に集中した測度であるから，
  $\delta_{T(0)} = \frac{1}{2}\delta_{-1} + \frac{1}{2}\delta_1$
  を満たす $T(0)$ は存在しない．

  \begin{center}
  \begin{tikzpicture}[scale=1.0]
    \draw[->] (-1.8, 0) -- (1.8, 0);
    \draw[thick, blue] (0, 0) -- (0, 1.2);
    \fill[blue] (0, 1.2) circle (2pt);
    \node[below] at (0, -0.1) {$0$};
    \node[above, blue] at (0, 1.35) {$\alpha = \delta_0$};

    \draw[->, thick, dashed, gray] (2.2, 0.6) -- (3.2, 0.6);
    \node[above, gray] at (2.7, 0.6) {$T\, ?$};

    \begin{scope}[xshift=5cm]
      \draw[->] (-1.8, 0) -- (1.8, 0);
      \draw[thick, red] (-1, 0) -- (-1, 0.6);
      \fill[red] (-1, 0.6) circle (2pt);
      \draw[thick, red] (1, 0) -- (1, 0.6);
      \fill[red] (1, 0.6) circle (2pt);
      \node[below] at (-1, -0.1) {$-1$};
      \node[below] at (1, -0.1) {$1$};
      \node[above, red] at (0, 0.8) {$\beta = \tfrac{1}{2}\delta_{-1} + \tfrac{1}{2}\delta_1$};
    \end{scope}
  \end{tikzpicture}
  \end{center}

  1点の質量を2箇所に分割して送ることができないため，
  実行可能な写像 $T$ 自体が存在しない．
\end{example}

この困難を解消するのが，次節の Kantorovich による緩和である．

%% ============================================================
%%  2. Kantorovich 緩和
%% ============================================================
\section{Kantorovich 緩和}
\label{sec:sem-kantorovich}

\subsection{鍵となるアイデア}

前節で見たとおり，Monge の定式化には
実行可能集合の非凸性と解の非存在という2つの本質的困難がある．
また，最適割当問題は $n = m$ かつ一様な重みの場合しか扱えず，
Monge 問題はこれを一般化するものの，$n!$ 通りの組合せ的構造を持つ．

Kantorovich (1942) は，これらの困難を一挙に解消するアイデアを提示した．
Monge の枠組みでは，各点 $x_i$ の質量をすべて1つの行き先 $T(x_i)$ に送る
\textbf{決定論的}な輸送を考えていた．
Kantorovich はこれを，質量を複数の行き先に分割して送ることを許す
\textbf{確率的}な輸送に置き換える．
この「\textbf{質量の分割}」（mass splitting）を許す輸送計画を，
$\X \times \Y$ 上の結合測度（\textbf{カップリング}）として定式化する．

\subsection{カップリング}

\begin{definition}{カップリング（離散版）}{sem-coupling-discrete}
  ヒストグラム $\mathbf{a} \in \simplex_n$，$\mathbf{b} \in \simplex_m$ に対して，
  $\mathbf{a}$ と $\mathbf{b}$ の\textbf{カップリング}の集合を
  \[
    \CouplingsD(\mathbf{a}, \mathbf{b}) \defeq
    \left\{ \mathbf{P} \in \R_+^{n \times m} \;\middle|\;
    \mathbf{P} \ones_m = \mathbf{a}, \;
    \mathbf{P}^\top \ones_n = \mathbf{b} \right\}
  \]
  で定義する．
  $P_{i,j}$ は「輸送元のビン $i$ から輸送先のビン $j$ へ送られる質量」を表す．
\end{definition}

行和の制約 $\mathbf{P} \ones_m = \mathbf{a}$ は
ビン $i$ から出る質量の合計が $a_i$ であることを，
列和の制約 $\mathbf{P}^\top \ones_n = \mathbf{b}$ は
ビン $j$ に届く質量の合計が $b_j$ であることを意味する．
$\CouplingsD(\mathbf{a}, \mathbf{b})$ は $n + m$ 本の等式制約で定まる
$\R_+^{n \times m}$ の有界部分集合であるから，
\textbf{凸多面体}（convex polytope）をなす．

離散版を一般の確率測度に拡張すると，次の連続版が得られる．

\begin{definition}{カップリング（連続版）}{sem-coupling}
  確率測度 $\alpha \in \Mm_+^1(\X)$ と $\beta \in \Mm_+^1(\Y)$ に対して，
  $\alpha$ と $\beta$ の\textbf{カップリング}の集合を
  \[
    \Couplings(\alpha, \beta) \defeq
    \left\{ \pi \in \Mm_+^1(\X \times \Y) \;\middle|\;
    (P_\X)\pushforward \pi = \alpha, \;
    (P_\Y)\pushforward \pi = \beta \right\}
  \]
  で定義する．ここで $P_\X \colon \X \times \Y \to \X$，
  $P_\Y \colon \X \times \Y \to \Y$ はそれぞれ第1・第2成分への射影である．
\end{definition}

すなわち，$\pi$ の周辺分布がそれぞれ $\alpha$, $\beta$ に一致することを要求する．

\begin{example}{Monge 写像が存在しない場合の解消}{sem-kantorovich-resolve}
  例~\ref{ex:sem-monge-noexist} では
  $\alpha = \delta_0$，$\beta = \frac{1}{2}\delta_{-1} + \frac{1}{2}\delta_1$
  に対して Monge 写像が存在しないことを見た．
  Kantorovich の枠組みでこの問題を考えよう．

  離散表現は $\mathbf{a} = (1)$，$\mathbf{b} = (1/2,\; 1/2)$ であり，
  カップリング $\mathbf{P} \in \R_+^{1 \times 2}$ は
  $\mathbf{P}\ones_2 = 1$，$\mathbf{P}^\top \ones_1 = \mathbf{b}$
  を満たす必要がある．行が1つしかないため $P_{1,j} = b_j$ が強制され，
  \[
    \mathbf{P} = \begin{pmatrix} 1/2 & 1/2 \end{pmatrix}
  \]
  が唯一のカップリングである．
  コスト $c(x,y) = |x - y|$ に対して
  \[
    \inner{\mathbf{C}}{\mathbf{P}}
    = \tfrac{1}{2} |0 - (-1)| + \tfrac{1}{2} |0 - 1| = 1.
  \]
  Monge の枠組みでは不可能だった「1点の質量を2箇所に分割する」操作が，
  Kantorovich のカップリングとして自然に実現されている．
\end{example}

\paragraph{Monge 写像との対比．}
\begin{figure}[htbp]
\centering
\begin{tikzpicture}[scale=0.95]
  % Monge (left)
  \node[above] at (1.5, 3.3) {\textbf{Monge 写像}};
  \fill[blue] (0.5, 2.5) circle (3pt) node[left] {\small $x_1$};
  \fill[blue] (0.5, 1.5) circle (3pt) node[left] {\small $x_2$};
  \fill[blue] (0.5, 0.5) circle (3pt) node[left] {\small $x_3$};
  \fill[red] (2.5, 2.5) circle (3pt) node[right] {\small $y_1$};
  \fill[red] (2.5, 1.5) circle (3pt) node[right] {\small $y_2$};
  \draw[->, thick] (0.7, 2.5) -- (2.3, 2.5);
  \draw[->, thick] (0.7, 1.5) -- (2.3, 2.5);
  \draw[->, thick] (0.7, 0.5) -- (2.3, 1.5);
  \node[below, text width=3cm, align=center] at (1.5, -0.2) {\small 各点は1つの\\行き先のみ};

  % Kantorovich (right)
  \begin{scope}[xshift=5.5cm]
    \node[above] at (1.5, 3.3) {\textbf{Kantorovich カップリング}};
    \fill[blue] (0.5, 2.5) circle (3pt) node[left] {\small $x_1$};
    \fill[blue] (0.5, 1.5) circle (3pt) node[left] {\small $x_2$};
    \fill[blue] (0.5, 0.5) circle (3pt) node[left] {\small $x_3$};
    \fill[red] (2.5, 2.5) circle (3pt) node[right] {\small $y_1$};
    \fill[red] (2.5, 1.5) circle (3pt) node[right] {\small $y_2$};
    \draw[->, thick] (0.7, 2.5) -- (2.3, 2.5);
    \draw[->, thick, blue!60] (0.7, 1.5) -- (2.3, 2.5);
    \draw[->, thick, red!60] (0.7, 1.5) -- (2.3, 1.5);
    \draw[->, thick] (0.7, 0.5) -- (2.3, 1.5);
    \node[right, font=\small] at (1.6, 1.8) {分割};
    \node[below, text width=3cm, align=center] at (1.5, -0.2) {\small 質量の分割を\\許す};
  \end{scope}
\end{tikzpicture}
\caption{Monge 写像と Kantorovich カップリングの対比．Monge 写像は各点を1つの行き先に送るが，Kantorovich カップリングは質量の分割を許す．}
\label{fig:sem-monge-vs-kantorovich}
\end{figure}

\subsection{Kantorovich 問題}

\begin{definition}{Kantorovich 問題（離散版）}{sem-kantorovich-discrete}
  ヒストグラム $\mathbf{a} \in \simplex_n$，$\mathbf{b} \in \simplex_m$，
  コスト行列 $\mathbf{C} \in \R_+^{n \times m}$ に対して，
  \textbf{離散 Kantorovich 問題}とは
  \[
    \MKD_{\mathbf{C}}(\mathbf{a}, \mathbf{b}) \defeq
    \min_{\mathbf{P} \in \CouplingsD(\mathbf{a}, \mathbf{b})}
    \inner{\mathbf{C}}{\mathbf{P}}
  \]
  を求める問題である．
  これは $nm$ 個の変数と $n + m$ 本の等式制約を持つ\textbf{線形計画問題}（LP）である．
\end{definition}

離散版を一般の確率測度に拡張すると，次の連続版が得られる．

\begin{definition}{Kantorovich 問題（連続版）}{sem-kantorovich}
  確率測度 $\alpha \in \Mm_+^1(\X)$，$\beta \in \Mm_+^1(\Y)$，
  コスト関数 $c \colon \X \times \Y \to \R_+$ に対して，
  \textbf{Kantorovich 問題}とは
  \[
    \MK_c(\alpha, \beta) \defeq
    \inf_{\pi \in \Couplings(\alpha, \beta)}
    \int_{\X \times \Y} c(x, y) \, \d\pi(x, y)
  \]
  を求める問題である．
\end{definition}

\subsection{最適割当問題との関係}

置換行列の定義から
$(\mathbf{P}_\sigma)_{i,j} = 1/n$ となるのは $j = \sigma(i)$ のときのみなので，
\[
  \inner{\mathbf{C}}{\mathbf{P}_\sigma}
  = \sum_{i,j} C_{i,j} (\mathbf{P}_\sigma)_{i,j}
  = \frac{1}{n} \sum_{i=1}^{n} C_{i,\sigma(i)}
\]
が成り立つ．すなわち最適割当問題は，カップリング $\mathbf{P}$ を
置換行列に制限した Kantorovich 問題に他ならない：
\[
  \min_{\sigma \in \Perm(n)} \frac{1}{n} \sum_{i=1}^{n} C_{i,\sigma(i)}
  = \min_{\sigma \in \Perm(n)} \inner{\mathbf{C}}{\mathbf{P}_\sigma}.
\]
置換行列の集合は $\CouplingsD(\ones_n/n, \ones_n/n)$ に含まれるから，
\[
  \MKD_{\mathbf{C}}(\ones_n/n, \ones_n/n)
  \leq \min_{\sigma \in \Perm(n)} \inner{\mathbf{C}}{\mathbf{P}_\sigma}
\]
である．次の命題は，この不等式で等号が成立することを示す．

\begin{proposition}{Kantorovich 緩和のタイト性}{sem-kantorovich-tight}
  $m = n$ かつ $\mathbf{a} = \mathbf{b} = \ones_n / n$ のとき，
  Kantorovich 問題の最適解として
  ある最適置換 $\sigma^* \in \Perm(n)$ に対応する
  置換行列 $\mathbf{P}_{\sigma^*}$ が存在する．
  すなわち，Kantorovich 緩和は最適割当問題に対して\textbf{タイト}である．
\end{proposition}

\begin{proof}
  Birkhoff の定理により，$\CouplingsD(\ones_n/n, \ones_n/n)$ の
  端点の集合は置換行列の全体に等しい．
  線形計画の基本定理により，空でない多面体上の線形目的関数の最小値は
  端点で達成される．
  したがって $\MKD_{\mathbf{C}}(\ones_n/n, \ones_n/n)$ の最適解として
  置換行列が取れる．
\end{proof}

\subsection{Kantorovich 緩和の利点}

Kantorovich の定式化は Monge 問題の困難をすべて解消する．

\begin{remark}{カップリング集合の凸性}{sem-coupling-convex}
  Monge の実行可能集合 $\mathcal{T}(\alpha, \beta)$ は非凸であったが，
  カップリング集合 $\CouplingsD(\mathbf{a}, \mathbf{b})$ は\textbf{凸}である．
  実際，$\mathbf{P}_1, \mathbf{P}_2 \in \CouplingsD(\mathbf{a}, \mathbf{b})$
  と $t \in [0,1]$ に対し，
  凸結合 $\mathbf{P} = t\mathbf{P}_1 + (1-t)\mathbf{P}_2$ は
  \[
    \mathbf{P}\ones_m = t\,\mathbf{a} + (1-t)\,\mathbf{a} = \mathbf{a}, \quad
    \mathbf{P}^\top \ones_n = t\,\mathbf{b} + (1-t)\,\mathbf{b} = \mathbf{b}
  \]
  を満たすから $\mathbf{P} \in \CouplingsD(\mathbf{a}, \mathbf{b})$ である．
  さらに目的関数 $\inner{\mathbf{C}}{\mathbf{P}}$ は $\mathbf{P}$ に関して線形であるから，
  離散 Kantorovich 問題は\textbf{凸最適化問題}（実際には線形計画問題）となる．
\end{remark}

解の存在も保証される．
$\CouplingsD(\mathbf{a}, \mathbf{b})$ は $\R^{n \times m}$ の有界閉集合であるから
コンパクトであり，連続関数の最小値が達成される．
さらに，積 $\mathbf{a}\mathbf{b}^\top$ は
$(\mathbf{a}\mathbf{b}^\top)\ones_m
= \mathbf{a}(\mathbf{b}^\top \ones_m) = \mathbf{a}$，
$(\mathbf{a}\mathbf{b}^\top)^\top \ones_n
= \mathbf{b}(\mathbf{a}^\top \ones_n) = \mathbf{b}$
を満たすためカップリングの1つを与え，実行可能集合は空でない．
したがって最適解 $\mathbf{P}^*$ が\textbf{常に存在}する．

\subsection{Monge 問題の緩和としての位置づけ}

Monge 写像 $T$ に対して，
$(\operatorname{Id}, T) \colon \X \to \X \times \Y$，$x \mapsto (x, T(x))$
により $\pi_T = (\operatorname{Id}, T)\pushforward \alpha$ とおくと
$\pi_T \in \Couplings(\alpha, \beta)$ となり，
\[
  \MK_c(\alpha, \beta) \leq
  \inf_{T}
  \left\{
  \int_\X c(x, T(x)) \, \d\alpha(x)
  \;\middle|\;
  T\pushforward \alpha = \beta
  \right\}.
\]
Kantorovich 問題は Monge 問題の緩和であり，
Monge 写像による輸送は「$\pi$ が $T$ のグラフ上に集中している」カップリングの特殊な場合である．

%% ============================================================
%%  3. Wasserstein 距離
%% ============================================================
\section{Wasserstein 距離}
\label{sec:sem-wasserstein}

Kantorovich 問題の最適値 $\MKD_{\mathbf{C}}(\mathbf{a}, \mathbf{b})$ は
ヒストグラム間の「輸送コスト」を与えるが，
コスト行列 $\mathbf{C}$ が距離行列のべき乗であるとき，
この値から自然に\textbf{距離の公理を満たす量}が定義できる．
すなわち，OT は点間の ground 距離を，分布間の距離に\textbf{持ち上げる}
（lift する）正準的な方法を提供する．

\subsection{Ground 距離}

\begin{definition}{Ground 距離行列}{sem-ground-metric}
  $\distD \in \R_+^{n \times n}$ が $\range{n}$ 上の\textbf{距離行列}であるとは，
  以下の3条件を満たすことをいう：
  \begin{enumerate}
    \item[(i)] 対称性：$\distD_{i,j} = \distD_{j,i}$，
    \item[(ii)] 正定値性：$\distD_{i,j} = 0 \iff i = j$，
    \item[(iii)] 三角不等式：$\distD_{i,k} \leq \distD_{i,j} + \distD_{j,k}$
      \quad $(\forall\, i, j, k \in \range{n})$．
  \end{enumerate}
\end{definition}

この距離行列 $\distD$ を「ビン間の距離」（ground metric）と呼ぶ．
例えば $\X = \R$ 上に台を持つヒストグラムでは，
$\distD_{i,j} = |x_i - x_j|$ のように台の点間のユークリッド距離を取ることが多い．

\subsection{$p$-Wasserstein 距離の定義}

\begin{proposition}{$p$-Wasserstein 距離}{sem-wasserstein-metric}
  $n = m$ とし，$p \geq 1$ に対して
  コスト行列を $\mathbf{C} = \distD^p = (\distD_{i,j}^p)_{i,j}$ とする．
  このとき
  \[
    \WassD_p(\mathbf{a}, \mathbf{b})
    \defeq \MKD_{\distD^p}(\mathbf{a}, \mathbf{b})^{1/p}
  \]
  は $\simplex_n$ 上の\textbf{距離}を定義する．
  すなわち，$\WassD_p$ は対称かつ正定値であり，三角不等式
  \[
    \WassD_p(\mathbf{a}, \mathbf{c})
    \leq \WassD_p(\mathbf{a}, \mathbf{b})
    + \WassD_p(\mathbf{b}, \mathbf{c})
    \qquad (\forall\, \mathbf{a}, \mathbf{b}, \mathbf{c} \in \simplex_n)
  \]
  を満たす．
\end{proposition}

\subsection{距離の公理の証明}

\begin{proof}
  \textbf{対称性と正定値性．}
  $\mathbf{C} = \distD^p$ の対角成分は $\distD_{i,i}^p = 0$ であるから，
  $\mathbf{P}^\star = \diag(\mathbf{a})$ が $\CouplingsD(\mathbf{a}, \mathbf{a})$ の
  実行可能解でコスト $0$ を達成し，$\WassD_p(\mathbf{a}, \mathbf{a}) = 0$．
  逆に $\mathbf{a} \neq \mathbf{b}$ のとき，
  ある添字 $i_0$ で $a_{i_0} \neq b_{i_0}$ であるから，
  周辺制約 $\mathbf{P}\ones = \mathbf{a}$, $\mathbf{P}^\top\ones = \mathbf{b}$ を満たす
  任意の $\mathbf{P}$ は対角成分だけでは $\mathbf{a} = \mathbf{b}$ を要求するため，
  必ず対角の外（$i \neq j$）に正の成分を持つ．
  $\distD^p$ の非対角成分はすべて正であるから $\WassD_p(\mathbf{a}, \mathbf{b}) > 0$．
  $\distD^p$ の対称性より $\WassD_p(\mathbf{a}, \mathbf{b}) = \WassD_p(\mathbf{b}, \mathbf{a})$．

  \medskip
  \textbf{三角不等式（glued coupling の構成）．}
  $\mathbf{a}, \mathbf{b}, \mathbf{c} \in \simplex_n$ とする．
  $\mathbf{P}$ を $(\mathbf{a}, \mathbf{b})$ 間の最適輸送，
  $\mathbf{Q}$ を $(\mathbf{b}, \mathbf{c})$ 間の最適輸送とする．

  \textbf{補助カップリングの構成．}
  $\tilde{\mathbf{b}}$ を $\tilde{b}_j = b_j$ ($b_j > 0$ のとき),
  $\tilde{b}_j = 1$ ($b_j = 0$ のとき) として，
  \[
    \mathbf{S} \defeq \mathbf{P} \, \diag(1 / \tilde{\mathbf{b}}) \, \mathbf{Q}
    \in \R_+^{n \times n}
  \]
  とおく．
  $\mathbf{S}$ は $(\mathbf{a}, \mathbf{c})$ のカップリングである：
  \[
    \mathbf{S} \ones_n
    = \mathbf{P} \, \diag(1/\tilde{\mathbf{b}}) \, \mathbf{Q} \ones_n
    = \mathbf{P} (\mathbf{b} / \tilde{\mathbf{b}})
    = \mathbf{P} \ones_{\mathrm{Supp}(\mathbf{b})}
    = \mathbf{a},
  \]
  ここで $\ones_{\mathrm{Supp}(\mathbf{b})}$ は
  $b_j > 0$ の添字で $1$，それ以外で $0$ のベクトルであり，
  $b_j = 0$ ならば周辺制約 $\sum_i P_{i,j} = b_j = 0$ と
  $P_{i,j} \geq 0$ から $P_{i,j} = 0$（$\forall\, i$）であるから
  $\mathbf{P} \ones_{\mathrm{Supp}(\mathbf{b})} = \mathbf{P} \ones = \mathbf{a}$
  が成り立つ．
  同様に $\mathbf{S}^\top \ones_n = \mathbf{c}$ も確認できる．

  \medskip
  \textbf{不等式の導出．}
  $\mathbf{S}$ は $\CouplingsD(\mathbf{a}, \mathbf{c})$ の実行可能解であるから（最適とは限らない），
  \begin{align*}
    \WassD_p(\mathbf{a}, \mathbf{c})
    &= \Bigl(\min_{\mathbf{P}' \in \CouplingsD(\mathbf{a}, \mathbf{c})}
       \inner{\mathbf{P}'}{\distD^p}\Bigr)^{1/p}
    \leq \inner{\mathbf{S}}{\distD^p}^{1/p} \\
    &= \Bigl(\sum_{i,k} \distD_{i,k}^p
       \sum_j \frac{P_{i,j} Q_{j,k}}{\tilde{b}_j}\Bigr)^{1/p}
    \leq \Bigl(\sum_{i,j,k}
       (\distD_{i,j} + \distD_{j,k})^p
       \frac{P_{i,j} Q_{j,k}}{\tilde{b}_j}\Bigr)^{1/p}
  \end{align*}
  ここで第1の不等式は $\mathbf{S}$ の次善性，
  第2の不等式は $\distD$ の三角不等式 $\distD_{i,k} \leq \distD_{i,j} + \distD_{j,k}$
  による．
  さらに Minkowski の不等式
  $\bigl(\sum_i |a_i + b_i|^p\bigr)^{1/p}
  \leq \bigl(\sum_i |a_i|^p\bigr)^{1/p} + \bigl(\sum_i |b_i|^p\bigr)^{1/p}$
  （$p \geq 1$）を，
  重み $P_{i,j} Q_{j,k} / \tilde{b}_j$ に関する加重和に適用すると
  \[
    \leq \Bigl(\sum_{i,j,k}
       \distD_{i,j}^p \frac{P_{i,j} Q_{j,k}}{\tilde{b}_j}\Bigr)^{1/p}
    + \Bigl(\sum_{i,j,k}
       \distD_{j,k}^p \frac{P_{i,j} Q_{j,k}}{\tilde{b}_j}\Bigr)^{1/p}.
  \]
  第1項では $k$ について和を取ると
  $\sum_k Q_{j,k} / \tilde{b}_j = b_j / \tilde{b}_j = 1$（$b_j > 0$ のとき），
  第2項では $i$ について和を取ると
  $\sum_i P_{i,j} / \tilde{b}_j = b_j / \tilde{b}_j = 1$ となり，
  \[
    \WassD_p(\mathbf{a}, \mathbf{c})
    \leq \Bigl(\sum_{i,j} \distD_{i,j}^p P_{i,j}\Bigr)^{1/p}
    + \Bigl(\sum_{j,k} \distD_{j,k}^p Q_{j,k}\Bigr)^{1/p}
    = \WassD_p(\mathbf{a}, \mathbf{b})
    + \WassD_p(\mathbf{b}, \mathbf{c}).
  \]
\end{proof}

証明の核心は，2つの最適輸送 $\mathbf{P}$, $\mathbf{Q}$ を
$\mathbf{b}$ を経由して「接着」（gluing）し，$(\mathbf{a}, \mathbf{c})$ 間の
実行可能な（しかし一般に最適ではない）カップリング $\mathbf{S}$ を構成する点にある．

\subsection{Wasserstein 距離の直観的理解}

$p$-Wasserstein 距離の重要な性質を確認する．

\paragraph{Dirac 測度間の距離．}
任意の2点 $x, y$ に対して
$\Couplings(\delta_x, \delta_y) = \{\delta_{(x,y)}\}$（1つしかない）であるから，
\[
  \Wass_p(\delta_x, \delta_y)^p = \dist(x, y)^p.
\]
すなわち $\Wass_p(\delta_x, \delta_y) = \dist(x,y)$ である．
点上の Dirac 測度間の Wasserstein 距離は，そのまま ground 距離に一致する．
これが「距離の持ち上げ」の最も基本的な例である．

\paragraph{弱距離としての性質．}
$L^2$ ノルムや KL ダイバージェンスなどの「古典的」な距離・ダイバージェンスは，
台が重ならない離散分布間では定義できないか，無限大となってしまう．
例えば $\delta_0$ と $\delta_\varepsilon$ はどちらも離散測度であり，
台が重ならないため $L^2$ ノルムは定義できず，KL ダイバージェンスは $+\infty$ となる．
一方，$\Wass_p(\delta_0, \delta_\varepsilon) = \varepsilon$ であるから，
$\varepsilon \to 0$ のとき $\Wass_p(\delta_0, \delta_\varepsilon) \to 0$ となり，
分布の「空間的なシフト」が自然に定量化される．
この性質が，OT を分布間の比較に用いる大きな動機である．

\begin{remark}{弱収束との関係}{sem-weak-convergence}
  コンパクト空間 $\X$ 上の確率測度の列 $(\alpha_k)_k$ が $\alpha$ に
  \textbf{弱収束}するとは，
  任意の連続関数 $g \in \Cc(\X)$ に対して
  $\int_\X g \, \d\alpha_k \to \int_\X g \, \d\alpha$
  が成り立つことをいう（$\alpha_k \rightharpoonup \alpha$ と記す）．
  コンパクト空間上では，$\Wass_p(\alpha_k, \alpha) \to 0$ と
  $\alpha_k \rightharpoonup \alpha$ は同値であることが知られている．
  すなわち Wasserstein 距離は弱収束を特徴づける距離である．
\end{remark}

\subsection{一般測度への拡張}

ヒストグラム上の結果は一般の測度に自然に拡張できる．

\begin{proposition}{測度上の $p$-Wasserstein 距離}{sem-wasserstein-measure}
  $\X = \Y$ とし，$\dist$ を $\X$ 上の距離，$c(x,y) = \dist(x,y)^p$ （$p \geq 1$）とする．
  このとき
  \[
    \Wass_p(\alpha, \beta)
    \defeq \MK_{\dist^p}(\alpha, \beta)^{1/p}
  \]
  は $\Mm_+^1(\X)$ 上の距離を定義する．
\end{proposition}

証明はヒストグラムの場合と同様に，
2つの最適カップリングを「接着」（gluing）して
中間測度を経由するカップリングを構成する手法に帰着する．
非コンパクト空間では，積分 $\int \dist(x,y)^p \, \d\pi$ が有限であることを
保証するため，$p$ 次モーメントが有限な測度，
すなわち
$\int_\X \dist(x, x_0)^p \, \d\alpha(x) < \infty$
（$x_0 \in \X$ は任意の基準点）を満たす $\alpha$ に制限して議論する．
